\section{多维变量变换与 Jacobian}
\label{sec:05-Probability/05-Joint-Distributions/05}

平面上有两个相互独立的标准正态坐标，比如一次定位的东西向误差和南北向误差。问：落点到原点的距离不超过一个标准差的概率是多少？

一维的记忆很牢：\(P(\abs{Z}\le 1)\approx 68\%\)。于是脱口而出的答案往往也是 \(68\%\)。实际是 \(39\%\)，差了将近一半。

错的不是概率，是几何。半径每往外挪一点，那一圈的周长就更长，能容纳的落点更多。密度沿半径确实在衰减，可面积在增长，两者相乘之后的分布形状与一维完全不是一回事。把这件事算准，需要一条规则：换一套坐标之后，密度该乘上什么。

\subsection{一维的答案是长度的缩放比}

一维的公式想必用过：\(Y=g(X)\)，反解 \(X=h(Y)\)，则 \(f_Y(y)=f_X(h(y))\,\abs{h'(y)}\)。它的来历是概率守恒——落进一小段区间的概率，在两套坐标下必须是同一个数：

\[
f_Y(y)\,dy=f_X(x)\,dx,\qquad dx=\abs{h'(y)}\,dy .
\]

导数在这里的角色不是``变化率''，而是局部的长度缩放比：\(y\) 那边一段长 \(dy\) 的区间，映回 \(x\) 那边长度变成 \(\abs{h'(y)}\,dy\)。密度是``概率除以长度''，长度被拉长了多少倍，密度就被稀释多少倍。

多维要守恒的东西一样，麻烦在于``长度''换成了``体积''，而体积的缩放不再由一个数决定。\(y\) 空间里一个小长方体映回 \(x\) 空间之后，一般不再是长方体，而是一个被扭斜的平行体。它的体积怎么算？

\subsection{小方块映成平行体，体积因子是行列式}

把逆变换 \(x=h(y)\) 在一点 \(y_0\) 附近做一阶展开，\(h(y_0+\Delta y)\approx h(y_0)+J_h(y_0)\Delta y\)，其中 \(J_h\) 是各分量对各 \(y_j\) 的偏导排成的矩阵。在足够小的尺度上，非线性映射就是一个线性映射——这正是可微的含义。

于是问题变成纯线性代数：以 \(\Delta y_1,\ldots,\Delta y_d\) 为边的小长方体，被矩阵 \(J_h\) 送成什么？答案是以 \(J_h\) 的各列（各乘上相应边长）为棱的平行体，而平行体的体积等于矩阵行列式的绝对值乘上原体积。

\begin{center}
\begin{tikzpicture}[x=1cm,y=1cm,font=\footnotesize,>=stealth]
  \draw[black!60,fill=blue!14] (0,0) rectangle (1.3,1.3);
  \draw[->,black!80,thick] (0,0) -- (1.3,0);
  \draw[->,black!80,thick] (0,0) -- (0,1.3);
  \node[black!75] at (0.65,-0.45) {体积 \(=\Delta y_1\Delta y_2\)};
  \draw[->,black!65] (2.0,0.65) -- (3.4,0.65);
  \node[black!75,above] at (2.7,0.72) {\(J_h\)};
  \begin{scope}[shift={(4.2,0)}]
    \draw[black!60,fill=orange!16] (0,0) -- (1.9,0.55) -- (2.4,1.95) -- (0.5,1.4) -- cycle;
    \draw[->,black!80,thick] (0,0) -- (1.9,0.55);
    \draw[->,black!80,thick] (0,0) -- (0.5,1.4);
    \node[black!75] at (1.3,-0.45) {体积 \(=\abs{\det J_h}\,\Delta y_1\Delta y_2\)};
  \end{scope}
\end{tikzpicture}

\emph{两条棱是 \(J_h\) 的两列。行列式在\hyperref[sec:03-Linear-Algebra/04-Determinants/01]{``行列式与体积：线性变换是放大、翻转，还是压扁空间''}一节被定义成有向体积的缩放因子，这里原样拿来用；取绝对值是因为密度不关心定向。}
\end{center}

把概率守恒写在这个小块上，两边约掉体积元并取极限，就得到多维密度变换公式。

\begin{theorem}
设随机向量 \(X\in\R^d\) 有密度 \(f_X\)，变换 \(Y=g(X)\) 在 \(X\) 的支撑上一一可逆、连续可微，且逆变换 \(h=g^{-1}\) 的 Jacobian 行列式处处非零。则 \(Y\) 的密度为
\[
f_Y(y)=f_X\bigl(h(y)\bigr)\,
\abs{\det \frac{\partial h}{\partial y}} .
\]
\end{theorem}

三个条件各挡一种失效。一一可逆挡的是``几个 \(x\) 挤到同一个 \(y\)''，这时该点的密度是多份贡献之和，单个 \(h\) 写不下；连续可微挡的是``局部无法用线性映射近似''，那样平行体的图像就不成立；行列式非零挡的是``体积被压成零''——变换把 \(d\) 维压进了更低维，像分布集中在一张曲面上，根本没有 \(\R^d\) 上的密度可言。三条都是为了让那张小方块的图能画出来。

\(d=1\) 时矩阵是一乘一的，行列式就是导数本身，一维公式原样浮现。而如果把这条定理里的密度换成一般的被积函数，它就是多重积分的换元公式（\hyperref[sec:02-Calculus/06-Multivariable-Calculus/07]{``多重积分与变量替换''}）。密度变换、换元积分、行列式作为体积因子，这三件事在教材里分属三章，内容是同一条。

\subsection{那个多出来的 \texorpdfstring{\(r\)}{r}}

回到开头。令 \(X=R\cos\Theta,\ Y=R\sin\Theta\)，从 \((r,\theta)\) 到 \((x,y)\) 的 Jacobian 行列式是 \(r\cos^2\theta+r\sin^2\theta=r\)，于是

\[
f_{R,\Theta}(r,\theta)=f_{X,Y}(r\cos\theta,\,r\sin\theta)\cdot r .
\]

\begin{center}
\begin{tikzpicture}[x=1cm,y=1cm,font=\footnotesize,>=stealth]
  \draw[->,black!55] (-0.2,0) -- (3.3,0) node[right] {\(r\)};
  \draw[->,black!55] (0,-0.2) -- (0,1.7) node[above] {\(\theta\)};
  \draw[black!60,fill=blue!25] (0.5,0.6) rectangle (0.9,1.0);
  \draw[black!60,fill=blue!25] (2.2,0.6) rectangle (2.6,1.0);
  \node[black!75] at (1.55,-0.62) {两块全等};
  \draw[->,black!65] (3.7,0.8) -- (4.7,0.8);
  \begin{scope}[shift={(5.6,0)}]
    \draw[->,black!55] (-0.2,0) -- (3.1,0) node[right] {\(x\)};
    \draw[->,black!55] (0,-0.2) -- (0,3.1) node[above] {\(y\)};
    \foreach \r in {0.5,0.9,2.2,2.6}{\draw[black!25] (2:\r) arc (2:88:\r);}
    \draw[black!25] (34:0.3) -- (34:2.9);
    \draw[black!25] (57:0.3) -- (57:2.9);
    \draw[black!60,fill=blue!25] (34:0.5) arc (34:57:0.5) -- (57:0.9) arc (57:34:0.9) -- cycle;
    \draw[black!60,fill=blue!25] (34:2.2) arc (34:57:2.2) -- (57:2.6) arc (57:34:2.6) -- cycle;
    \node[black!75] at (1.55,-0.62) {像的面积相差三倍以上};
  \end{scope}
\end{tikzpicture}

\emph{左边两块在 \((r,\theta)\) 坐标下完全一样大；映到平面上，外圈那块的面积约为 \(r\,\Delta r\,\Delta\theta\)，随 \(r\) 线性增长。这个增长率就是 \(\abs{\det J}=r\)。忘了它，等于默认内圈和外圈一样宽敞。}
\end{center}

对独立标准正态代入，联合密度 \(\frac{1}{2\pi}e^{-(x^2+y^2)/2}\) 变成

\[
f_{R,\Theta}(r,\theta)=\frac{1}{2\pi}\cdot r e^{-r^2/2},
\]

右边恰好分成只含 \(\theta\) 的因子与只含 \(r\) 的因子——按本单元第二节的因式分解判据，\(R\) 与 \(\Theta\) 独立，\(\Theta\) 在 \([0,2\pi)\) 上均匀，而 \(f_R(r)=re^{-r^2/2}\)。这个密度在 \(r=0\) 处等于零：原点虽然是联合密度最高的点，却几乎没有面积可分。积一下便得

\[
P(R\le 1)=\int_0^1 re^{-r^2/2}\,dr=1-e^{-1/2}\approx 0.39 ,
\]

开头那个反直觉的数字就是这么来的。它同时说明了一件更一般的事：维度升高时，按欧氏距离衡量的概率质量会往外挪，``离中心一个标准差以内''这类判断不能跨维度照搬。顺带一提，把这个分解反过来用——先抽均匀的角度、再按 \(f_R\) 抽半径——就是 Box--Muller 生成正态样本的做法。

\subsection{和变换：为什么卷积公式长那样}

本单元第三节的卷积公式当时是直接给出的，用这条定理可以补上来历。引入辅助变量，令 \(U=X+Y,\ V=X\)，逆变换是 \(x=v,\ y=u-v\)，Jacobian 行列式的绝对值等于 \(1\)。于是

\[
f_{U,V}(u,v)=f_{X,Y}(v,\,u-v),
\]

再把 \(V\) 边缘化掉就得到 \(f_U(u)=\int f_{X,Y}(v,u-v)\,dv\)，独立时进一步化成卷积。行列式等于 \(1\) 解释了为什么当初那条式子里没有任何多余因子——剪切变换不改变面积。

这里的套路比结果更值得记：想求某个函数 \(U=\psi(X,Y)\) 的分布，就补一个辅助变量凑成可逆变换，套公式得到联合密度，再把辅助变量积掉。比对 CDF 求导要机械得多，遇到商、比值、极大值之类的变换都能照走。

\subsection{在生成模型里，这个因子就是训练目标的一半}

现代生成模型把这条公式用到了明处。可逆流模型（normalizing flow）的做法是：从一个简单分布（通常是标准正态）取 \(z\)，经过一串可逆变换 \(x=h(z)\) 得到样本。训练要最大化真实数据的对数似然，而按上面的定理取对数，

\[
\log f_X(x)=\log f_Z\bigl(h^{-1}(x)\bigr)+\log\abs{\det \frac{\partial h^{-1}}{\partial x}} .
\]

第一项要求变换把数据拉回到正态的高密度区，第二项是体积记账——如果网络靠着把空间疯狂压缩来堆高密度，这一项会给出相应的惩罚。整个架构设计的难点也集中在这一项：既要让变换足够灵活，又要让 \(d\times d\) 的行列式算得动，于是才有了各种把 Jacobian 逼成三角阵的耦合层设计。

同一条公式的轻量版本更常见。重参数化 \(x=\mu+\sigma z\) 是最简单的可逆变换，它的 \(\log\abs{\det J}\) 就是 \(\sum_i\log\sigma_i\)——变分自编码器里那一项看似突兀的对数标准差之和，来历就在这里。

\subsection{不可逆时怎么办}

映射多对一时，公式要按分支相加：找出所有映到同一个 \(y\) 的逆像 \(h_1(y),h_2(y),\ldots\)，各算各的因子再求和，

\[
f_Y(y)=\sum_k f_X\bigl(h_k(y)\bigr)\,\abs{\det \frac{\partial h_k}{\partial y}} .
\]

一维的平方变换就是这个形式的最小例子。极坐标之所以不需要求和，是因为我们事先把 \(r\ge 0\)、\(\theta\in[0,2\pi)\) 限死了；放开限制，同一个 \((x,y)\) 会有无穷多种 \((r,\theta)\) 表示，公式随即失效。定义域的选择不在公式的职责范围内，那始终是使用者的事。

若变换根本不可逆，或者在某片区域上 Jacobian 退化，就得退回更原始的手段：直接算 \(P(Y\in A)\)，或把定义域切成若干块分别处理。这类情形并不罕见——把 \(d\) 维数据投影到二维、取分量的极大值、做四舍五入，都不可逆。此时``\(Y\) 的密度''可能压根不存在，硬套公式会得到一个看起来正常、实则无意义的表达式。

按顺序执行能避开大多数坑：写正变换、解逆变换、确定新变量的支撑、算行列式的绝对值、代入、最后用全空间积分是否等于 \(1\) 做一次核对。经验上出错最多的是支撑范围，其次是漏掉绝对值或忘了分支求和；求偏导本身反而最不容易错。

\medskip

本单元到此配齐了处理多个随机变量的五件工具：联合分布保存配对，独立性把高维拆成低维，卷积回答和的分布，协方差矩阵概括二阶几何，Jacobian 管住坐标变换下的密度。它们都是关于结构的，不依赖具体是哪一个分布。下一章反过来，去认那些在实际问题中反复出现的具体分布——先辨认背后的随机机制，再决定用哪一个。
